\chapter{The Methods of Experimental Design}
\label{ch:design}

\begin{quote}\itshape
``To consult the statistician after an experiment is finished is often merely to
ask him to conduct a post-mortem examination. He can perhaps say what the
experiment died of.'' --- R.~A.~Fisher.
\end{quote}

\noindent
A mechanistic model, however faithful, is silent until it is \emph{queried}. The
science of \emph{how} to query it---which combinations of process conditions to
evaluate, in what order, and how to turn the results into knowledge---is
\emph{Design of Experiments}. In a development setting the same model is the
oracle of Chapter~\ref{ch:models} or a real wet-lab run; the design methods are
identical. This chapter presents the families the package implements, in the
order one typically reaches for them: \emph{full factorial designs} with
response-surface analysis (Section~\ref{sec:factorial}) for screening and
robustness; \emph{generalized linear models} (Section~\ref{sec:glm}) for the
pass/fail and count responses that yield a \emph{probabilistic design space};
\emph{Latin hypercube sampling} (Section~\ref{sec:lhs}) for space-filling
exploration; \emph{latent-variable analysis}---PCA and PLS---(Section~\ref{sec:mv})
for understanding a high-dimensional design table; and \emph{Bayesian
optimisation} (Section~\ref{sec:bo}) for adaptively hunting an optimum.
Section~\ref{sec:design-compare} compares them.

A common vocabulary, borrowed from \emph{Quality by Design} (QbD), runs through
all four. The inputs we vary are the \emph{critical process parameters} (CPPs):
load pH, salt, gradient slope, load density. The outputs we measure are the
\emph{critical quality attributes} (CQAs) and performance metrics: yield, purity,
productivity. The region of CPP space within which the CQAs remain acceptable is
the \emph{design space}, and the project of process characterisation is to map
it. In the package a CPP is a \code{Factor}---a name and a physical low/high
bound---and every design is generated over a list of such factors.

%==============================================================================
\section{Full factorial designs and response-surface analysis}
\label{sec:factorial}

\subsection{Constructing the design}

A \emph{two-level full factorial} design evaluates every combination of the
extreme settings of each factor. With $k$ factors this is $2^k$ runs. Each factor
is expressed in \emph{coded} units that map its physical low and high to $-1$ and
$+1$; the affine decoding (method \code{Factor.decode}) is
%
\begin{equation}
x_{\mathrm{phys}}
= \underbrace{\tfrac{1}{2}(x_{\mathrm{hi}}+x_{\mathrm{lo}})}_{\text{centre } c_0}
+ x_{\mathrm{coded}}\cdot
\underbrace{\tfrac{1}{2}(x_{\mathrm{hi}}-x_{\mathrm{lo}})}_{\text{half-range } \Delta}.
\label{eq:decode}
\end{equation}
%
Coding is not cosmetic: it makes the design \emph{orthogonal}, so that the
estimated effects are statistically independent (their estimates do not change as
terms are added or dropped) and directly comparable in magnitude. The package's
\code{full\_factorial} enumerates the $2^k$ (or, at three levels, $3^k$) coded
points with \code{itertools.product} and decodes them. Optional replicated
\emph{centre points}, placed at every factor's midpoint, are appended; they serve
two purposes---they provide an estimate of \emph{pure experimental error} from
replication, and, by comparing the centre response to the factorial average, a
test for \emph{curvature}. Finally the run order is randomised (with a fixed
seed) so that any uncontrolled time trend is converted into noise rather than
bias. The geometry is the familiar cube of Figure~\ref{fig:cube}.

\begin{figure}[t]
\centering
\includegraphics[width=0.6\linewidth]{factorial_cube.png}
\caption[The $2^3$ factorial design]{A $2^3$ full-factorial design in three
coded factors: the eight vertices of the cube plus a replicated centre point.
Every main effect and every interaction is estimable, and the design is
orthogonal and rotatable about the centre.}
\label{fig:cube}
\end{figure}

\subsection{The response-surface model}

The responses are analysed (\code{doe/analysis.py}) by ordinary least squares,
fitting a polynomial \emph{response surface}. For the general second-order model,
%
\begin{equation}
y = \beta_0
  + \sum_{i} \beta_i x_i
  + \sum_{i<j} \beta_{ij} x_i x_j
  + \sum_i \beta_{ii} x_i^2
  + \varepsilon,
\qquad \varepsilon \sim \mathcal N(0,\sigma^2),
\label{eq:rsm}
\end{equation}
%
where the $\beta_i$ are \emph{main effects}, the $\beta_{ij}$ \emph{two-factor
interactions}, and the $\beta_{ii}$ \emph{quadratic} (curvature) terms. The model
is built programmatically in Wilkinson--Rogers notation and fitted with
\code{statsmodels}; main effects and interactions are included by default,
quadratics on request (the latter requiring a three-level or central-composite
design, since a two-level design cannot resolve curvature). In matrix form
$\bm y = \bm X\bm\beta + \bm\varepsilon$, the least-squares estimate and its
covariance are the classical
%
\begin{equation}
\hat{\bm\beta} = (\bm X\transp\bm X)^{-1}\bm X\transp \bm y,
\qquad
\Var(\hat{\bm\beta}) = \sigma^2 (\bm X\transp\bm X)^{-1},
\label{eq:ols}
\end{equation}
%
and orthogonal coding makes $\bm X\transp\bm X$ diagonal, so each coefficient is
estimated with the same, minimal variance.

\subsection{Apportioning variance: the ANOVA}

To judge which effects are real the package computes a \emph{Type-II analysis of
variance} (\code{anova\_lm(\dots, typ=2)}). Total variability in the response is
partitioned into a part explained by each term and an unexplained residual; the
significance of a term is judged by an $F$-ratio comparing its mean square to the
residual mean square,
%
\begin{equation}
F = \frac{\mathrm{SS}_{\text{effect}}/\mathrm{df}_{\text{effect}}}
         {\mathrm{SS}_{\text{residual}}/\mathrm{df}_{\text{residual}}}.
\end{equation}
%
The Type-II sum of squares measures each effect \emph{adjusted for all others of
equal or lower order}, which is the correct choice for the orthogonal factorial
designs used here and is standard in process-characterisation practice.

\subsection{Proven acceptable ranges}

The practical output of a robustness study is, for each factor, the range over
which it may drift while the response stays in specification---its \emph{proven
acceptable range} (PAR). The package (\code{proven\_acceptable\_ranges}) derives
a univariate PAR by inverting the fitted surface: holding all other factors at
their centre, it scans factor $x_i$ across its studied range and reports the
sub-range over which the prediction satisfies $\mathrm{spec}_{\mathrm{lo}} \le
\hat y(x_i) \le \mathrm{spec}_{\mathrm{hi}}$. It is a first-order, one-factor-at-a-
time reading of the model; a genuinely multivariate design space is read instead
from the response-surface contours.

\begin{prosbox}{factorial designs}
\begin{itemize}
  \item \textbf{Orthogonal and efficient.} Every main effect and interaction is
  estimated independently and with minimum variance; the analysis is transparent
  and the coefficients are directly interpretable.
  \item \textbf{Interactions, by construction.} Unlike one-factor-at-a-time
  experimentation, a factorial reveals how factors combine---often the most
  important finding of a study.
  \item \textbf{Built-in diagnostics.} Centre points give a model-free estimate
  of error and a curvature check; the ANOVA gives formal significance tests.
  \item \textbf{Regulatory pedigree.} Factorial and response-surface designs are
  the lingua franca of QbD submissions.
\end{itemize}
\end{prosbox}

\begin{consbox}{factorial designs}
\begin{itemize}
  \item \textbf{Exponential cost.} $2^k$ (or $3^k$) runs become prohibitive
  beyond a handful of factors; fractional factorials trade this against
  \emph{confounding} of effects, which a full design avoids only at full cost.
  \item \textbf{Low-order model.} A polynomial surface captures gentle curvature
  but not the sharp, threshold-like responses a mechanistic model can
  produce---it is a local approximation, valid inside the studied box.
  \item \textbf{Corners only (two-level).} With no interior points a two-level
  design cannot detect curvature at all; it must be augmented to a three-level or
  central-composite design.
  \item \textbf{Static.} The whole design is committed up front; it does not
  learn from early results.
\end{itemize}
\end{consbox}

%==============================================================================
\section{Generalized linear models and the probabilistic design space}
\label{sec:glm}

The response surface of Section~\ref{sec:factorial} models a \emph{continuous}
response---yield, log-purity, productivity---with Gaussian scatter. But the
question a regulatory filing ultimately asks is not ``what is the mean purity?''
but ``\emph{what is the probability that this run meets specification?}'' That
response is a \emph{yes/no}, and a second one---the number of impurity peaks that
contaminate the product pool---is a \emph{count}. Neither is continuous-and-normal,
and forcing a least-squares surface onto them is the error
Section~\ref{sec:stats-glm} warned against: a linear model of a $0/1$ outcome
predicts probabilities outside $[0,1]$, and one of a count predicts negative
counts. The right instrument is the \emph{generalized linear model} (GLM).

\subsection{From a response surface to a probability surface}

Recall the GLM of Section~\ref{sec:stats-glm}: the same linear predictor in the
process factors, $\eta(\bm x)=\beta_0+\sum_i\beta_i x_i+\sum_{i<j}\beta_{ij}x_ix_j$
(main effects, interactions, and---with a three-level design---curvature, exactly
as in the response surface), passed through a \emph{link} appropriate to the
response. For a pass/fail outcome the logit link gives a \emph{logistic
response surface},
%
\begin{equation}
P(\text{meet spec}\mid\bm x) = \sigma\!\bigl(\eta(\bm x)\bigr)
   = \frac{1}{1+e^{-\eta(\bm x)}},
\label{eq:glm-logit}
\end{equation}
%
a smooth surface that lives, by construction, in $[0,1]$. For the count of
co-eluting impurities the log link gives a \emph{Poisson response surface},
$\E[\text{count}\mid\bm x]=e^{\eta(\bm x)}$, automatically non-negative. Both are
fitted by maximum likelihood (the iteratively reweighted least squares of
Section~\ref{sec:stats-glm}) through the package's \code{fit\_glm\_response}
(\code{doe/analysis.py}), which builds the Wilkinson--Rogers formula exactly as
the OLS \code{fit\_response\_model} does and hands it to a \code{statsmodels} GLM
with the chosen family. Goodness of fit is read not from $R^2$ but from the
\emph{deviance} and McFadden's pseudo-$R^2$, $1-\ell/\ell_0$, the log-likelihood
relative to an intercept-only model.

Equation~\eqref{eq:glm-logit} is the modern, probabilistic reading of the
\emph{design space} of Quality by Design~\cite{ich}. Where the proven acceptable
ranges of Section~\ref{sec:factorial} carve out, deterministically, where the
\emph{mean} response sits in spec, the logistic surface delivers the
\emph{assurance} region
%
\begin{equation}
\mathrm{DS}_{1-\alpha}
  = \bigl\{\,\bm x : P(\text{meet all CQAs}\mid\bm x)\ \ge\ 1-\alpha\,\bigr\},
\label{eq:design-space}
\end{equation}
%
the set of process conditions at which the run passes with high probability (say
$1-\alpha=0.9$). This is precisely the deliverable a process-characterisation
study is meant to produce, and it falls straight out of a GLM fitted to pass/fail
data.

\subsection{A polishing-step application: purity}
\label{sec:glm-polish}

Consider a cation-exchange polish that must separate a monoclonal antibody from
three closely related \emph{charge variants}---the textbook hard case, since the
variants differ from the product by a whisker of surface charge. We expose two
critical process parameters: the \emph{load pH} (which, because the variants
carry slightly different pH-sensitivities, is the \emph{selectivity} lever) and
the \emph{gradient length} in column volumes (the \emph{throughput/yield} lever).
A Latin-hypercube design (Section~\ref{sec:lhs}) places $64$ runs across this
two-factor space; each run is a full \code{run\_column} simulation of the
four-component separation. From each chromatogram we pool a fixed window about
the target apex and record two responses: whether the pool clears the
specification ($\text{purity}\ge0.95$ \emph{and} $\text{yield}\ge0.80$), and how
many variants co-elute under the target peak.

Figure~\ref{fig:chrom-glm} shows the two GLMs. The logistic surface (left) maps a
crisp probabilistic design space: a band at \emph{low} pH---where the variants are
pushed apart and the pool is pure---bounded below by a yield cliff at the steepest
gradients, so that the $P\ge0.9$ contour (dashed) encloses a modest corner of the
operating window. The fit is excellent (pseudo-$R^2\approx0.86$), and the
interaction between pH and gradient is real: gradient length matters for the
verdict only at the low-pH edge, where yield is in play. The Poisson surface
(right) tells the mechanistic story behind the failures: the expected number of
co-eluting variants climbs steeply as pH rises and selectivity collapses
(pseudo-$R^2\approx0.59$). Together they convert a few dozen simulated runs into
an interpretable, probabilistic map of where the step can be operated---and why it
fails where it does.

\begin{keybox}{Algorithm: a probabilistic design space from pass/fail data}
\begin{enumerate}
  \item \textbf{Design \& simulate.} Lay an LHS (or factorial) over the CPPs;
  run the chromatography oracle at each point and reduce each chromatogram to its
  pool metrics.
  \item \textbf{Score.} Encode the binary verdict $y_k=\mathbf 1[\text{purity}_k
  \ge p^\star \wedge \text{yield}_k\ge y^\star]$ and, optionally, the impurity
  count.
  \item \textbf{Fit.} \code{fit\_glm\_response(df, "passed", factors,
  family="binomial")} (and \code{family="poisson"} for the count); maximum
  likelihood via IRLS returns the coefficients and the deviance.
  \item \textbf{Map.} \code{predict\_glm\_grid} evaluates
  Eq.~\eqref{eq:glm-logit} over a grid; thresholding at $1-\alpha$ gives the
  design space~\eqref{eq:design-space}.
\end{enumerate}
\end{keybox}

\begin{figure}[t]
\centering
\includegraphics[width=\linewidth]{chrom_glm.png}
\caption[GLM analysis of a chromatography design space]{Generalized linear models
turn $64$ simulated cation-exchange runs into a probabilistic design space.
\emph{Left:} a \emph{logistic} response surface for $P(\text{meet spec})$ over load
pH and gradient length; markers are the design runs (pass $\circ$, fail
$\times$), and the dashed line is the $P=0.9$ boundary of the design
space~\eqref{eq:design-space}. \emph{Right:} a \emph{Poisson} response surface for
the expected number of charge variants co-eluting under the target peak, which
rises sharply as rising pH erodes selectivity. Computed by \code{run\_column} and
fitted by \code{fit\_glm\_response}.}
\label{fig:chrom-glm}
\end{figure}

\subsection{A capture-step application: breakthrough and recovery}
\label{sec:glm-capture}

The polish step above is judged on \emph{purity}. The first chromatography step
of a process---the \emph{capture}---is judged on \emph{recovery}, and its
characteristic failure is different: \emph{breakthrough}. Run in bind-and-elute
mode, a capture column is loaded with feed until just before the product begins
to escape unbound at the outlet; load too much, and the column saturates and the
product is lost to the flow-through (Section~\ref{sec:sep-column}). The amount
that can be loaded before this happens---the \emph{dynamic binding
capacity}---is the quantity a capture step is designed around, and it depends on
how much is loaded and on how strongly the molecule binds.

We expose two critical process parameters of a cation-exchange capture: the
\emph{load density} (grams of product offered per litre of resin) and the
\emph{load conductivity}, set by the salt concentration. Conductivity is the
telling one: because ion-exchange binding weakens as salt rises
(Section~\ref{sec:iex}), a high-conductivity load erodes the binding strength and
the product breaks through early---which is exactly why capture loads are kept
low in salt. Figure~\ref{fig:capture-glm}(a) shows the three regimes as
\emph{breakthrough curves}---the outlet concentration, relative to feed, against
the volume loaded: a clean capture that never breaks through, an overload whose
front emerges once the bed saturates, and a weak-binding case that bleeds product
from early on.

For each run we record the \emph{recovery}---one minus the fraction lost to the
flow-through---and ask whether it clears a specification of $90\%$. Here a subtle
but important point arises, and it is the reason the design space is
\emph{probabilistic} at all. The simulated recovery is, in itself, a smooth
deterministic function of the two parameters; a single noiseless run at each
condition would partition the space into pass and fail by a sharp curve, and a
logistic fit to it would be degenerate (perfectly separable). Real runs are not
noiseless: replicate loads of nominally identical conditions scatter, through
assay error and the batch-to-batch variability of
Section~\ref{sec:stats-varcomp}. Drawing a handful of noisy replicate
measurements at each design point---the package's measurement-noise model
(\code{add\_measurement\_noise})---turns the verdict near the boundary into a
genuine probability, and the logistic GLM recovers it cleanly
(Figure~\ref{fig:capture-glm}(b), pseudo-$R^2\approx0.88$). The $P\ge0.9$ contour
is the \emph{assurance} boundary~\eqref{eq:design-space}: the curved frontier in
load-density/conductivity space inside which the capture step meets its recovery
target with high confidence. This is the probabilistic design space in its most
operationally honest form---its very gradient is a statement about run-to-run
variability, not just about the mean.

\begin{figure}[t]
\centering
\includegraphics[width=\linewidth]{capture_glm.png}
\caption[A capture-step breakthrough design space]{A cation-exchange capture step
analysed by logistic regression. \emph{Left:} representative \emph{breakthrough
curves}---outlet over feed concentration versus volume loaded---for a clean
capture, an overload, and a weak-binding (high-salt) load; the dotted line marks
$10\%$ breakthrough. \emph{Right:} the logistic \emph{recovery design space} over
load density and load conductivity, fitted to replicate noisy assays of $48$
simulated runs; filled contours are $P(\text{recovery}\ge90\%)$, markers are the
design points shaded by their observed pass fraction, and the dashed line is the
$P=0.9$ assurance boundary. Computed by \code{run\_column} and
\code{fit\_glm\_response}.}
\label{fig:capture-glm}
\end{figure}

\begin{prosbox}{generalized linear models}
\begin{itemize}
  \item \textbf{Right tool for the response.} Models pass/fail and count data on
  their own terms---bounded probabilities, non-negative counts---where a
  least-squares surface is simply invalid.
  \item \textbf{The probabilistic design space, directly.} A logistic surface is
  the assurance map~\eqref{eq:design-space} that QbD asks for, and it inherits the
  interaction and curvature structure of the response-surface model.
  \item \textbf{Interpretable and cheap.} Signed coefficients (log-odds, log-rate)
  with standard errors and a deviance test; fitted in milliseconds on top of the
  simulations.
\end{itemize}
\end{prosbox}

\begin{consbox}{generalized linear models}
\begin{itemize}
  \item \textbf{Needs events of both kinds.} A logistic fit is uninformative if
  every run passes (or every run fails); the design must straddle the boundary.
  \item \textbf{Still a low-order surface.} Like the OLS response surface it is a
  smooth local model; sharp, non-monotone boundaries need interactions, curvature,
  or more data.
  \item \textbf{Threshold-dependent.} The design space moves with the spec and the
  assurance level $1-\alpha$; it is a reading of the model, not an absolute.
\end{itemize}
\end{consbox}

%==============================================================================
\section{Space-filling: Latin hypercube sampling}
\label{sec:lhs}

When the goal is to \emph{explore} a response over its whole domain---to build a
surrogate, or to feed a global sensitivity analysis---rather than to fit a
low-order polynomial at the corners, one wants points spread as evenly as
possible through the interior. Placing them on a full grid is infeasible (the
curse of dimensionality: $\ell$ levels in $k$ dimensions is $\ell^k$ points), and
independent random sampling leaves clumps and gaps. \emph{Latin hypercube
sampling} (LHS) is the standard compromise.

\subsection{Construction}

An LHS of $n$ points in $k$ dimensions partitions each axis into $n$ equal
intervals and places exactly one point in each interval of each axis: a
\emph{Latin square} generalised to $k$ dimensions. The marginal projection onto
any single factor is therefore perfectly stratified---every one of the $n$ bins
is occupied---no matter how large $k$ is. Formally, with $\pi_1,\dots,\pi_k$
independent random permutations of $\{0,\dots,n-1\}$ and $U_{ij}\sim\mathcal
U[0,1)$,
%
\begin{equation}
x_{ij} = \frac{\pi_j(i) + U_{ij}}{n},
\end{equation}
%
scaled afterwards to the physical bounds (\code{qmc.scale}). Plain LHS guarantees
good \emph{one}-dimensional projections but can still correlate factors or leave
multi-dimensional voids; the package therefore uses
\code{scipy.stats.qmc.LatinHypercube} with the \code{"random-cd"} criterion,
which permutes the design to minimise the \emph{centered $L_2$ discrepancy}---a
global measure of departure from uniformity,
%
\begin{equation}
D_{\mathrm{CD}}^2 =
\left(\tfrac{13}{12}\right)^{k}
- \frac{2}{n}\sum_{i=1}^{n}\prod_{j=1}^{k}
  \Bigl(1 + \tfrac12|x_{ij}-\tfrac12| - \tfrac12|x_{ij}-\tfrac12|^2\Bigr)
+ \frac{1}{n^2}\sum_{i=1}^{n}\sum_{l=1}^{n}\prod_{j=1}^{k}
  \Bigl(1 + \tfrac12|x_{ij}-\tfrac12| + \tfrac12|x_{lj}-\tfrac12|
        - \tfrac12|x_{ij}-x_{lj}|\Bigr).
\label{eq:cd}
\end{equation}
%
Lower discrepancy means more uniform coverage in \emph{all} projections at once.

\subsection{Coverage diagnostics}

The package (\code{coverage\_metrics}) reports the centered discrepancy of
Eq.~\eqref{eq:cd} together with the minimum and mean pairwise Euclidean distance
between points (on factors normalised to the unit cube). A good space-filling
design has low discrepancy and a large minimum distance (no two points crowd
together). Figure~\ref{fig:lhs} contrasts an optimised LHS with an
equal-sized random sample: the random design clumps and gaps, while the LHS
covers the square evenly and achieves a markedly lower discrepancy.

\begin{figure}[t]
\centering
\includegraphics[width=0.92\linewidth]{lhs_vs_random.png}
\caption[LHS versus random sampling]{Twenty points in a two-factor space.
\emph{Left:} independent random sampling, which leaves visible clusters and
voids. \emph{Right:} an optimised Latin hypercube, which both stratifies every
marginal and achieves a lower centered discrepancy $D_{\mathrm{CD}}$. The
advantage widens as the dimension grows.}
\label{fig:lhs}
\end{figure}

\begin{prosbox}{Latin hypercube sampling}
\begin{itemize}
  \item \textbf{Dimension-robust uniformity.} Perfect one-dimensional
  stratification for any $n$ and $k$, with $n$ chosen freely (not tied to
  $\ell^k$).
  \item \textbf{Ideal surrogate-training data.} Even interior coverage is exactly
  what a Gaussian process or other global model needs.
  \item \textbf{Projection-friendly.} If some factors turn out inert, the design
  remains well-spread in the active subspace.
\end{itemize}
\end{prosbox}

\begin{consbox}{Latin hypercube sampling}
\begin{itemize}
  \item \textbf{No structured estimability.} It is not built to estimate specific
  polynomial effects; there is no clean ANOVA decomposition.
  \item \textbf{Optimisation needed for high dimensions.} Naive LHS can still
  correlate factors; the discrepancy-minimising step is essential and costs
  search time.
  \item \textbf{Not adaptive.} Like the factorial it is a one-shot design that
  does not concentrate effort where the response is interesting.
  \item \textbf{Edge under-sampling.} Points sit at bin centres, so the extreme
  corners of the domain---sometimes the most important---are sampled less than in
  a factorial.
\end{itemize}
\end{consbox}

%==============================================================================
\section{Latent-variable analysis: PCA and PLS}
\label{sec:mv}

Once a design table exists---factors and responses, possibly many of each, on
different scales and correlated---one needs tools to \emph{see} its structure and
to relate inputs to outputs. The package provides the two workhorses of
multivariate process analysis (\code{doe/multivariate.py}), both of which begin by
\emph{auto-scaling} every column to zero mean and unit variance so that variables
of different physical units contribute on equal footing.

\subsection{Principal-component analysis}

PCA finds the orthogonal directions of greatest variance in the (scaled) factor
matrix $\bm X$. It factorises
%
\begin{equation}
\bm X_s = \bm T \bm P\transp + \bm E,
\end{equation}
%
where the columns of the \emph{loadings} $\bm P$ are the eigenvectors of the
sample covariance $\tfrac{1}{n-1}\bm X_s\transp\bm X_s$, ordered by decreasing
eigenvalue, and the \emph{scores} $\bm T = \bm X_s\bm P$ are the coordinates of
each run in the new, decorrelated basis. The fraction of total variance carried
by component $a$ is $\lambda_a / \sum_b \lambda_b$ (the
\code{explained\_variance\_ratio}); the first few components usually capture most
of it, so a high-dimensional table can be \emph{visualised} in two or three score
dimensions, with clusters, trends, and outliers made plain.

\subsection{Partial least squares}

PCA looks only at $\bm X$. To \emph{relate} factors to responses one wants the
directions in $\bm X$ that best predict $\bm Y$, and that is PLS. It builds a
shared set of latent variables that decompose both blocks,
%
\begin{equation}
\bm X_s = \bm T \bm P\transp + \bm E,
\qquad
\bm Y_s = \bm T \bm Q\transp + \bm F,
\end{equation}
%
choosing, at each step, the weight vector $\bm w$ that maximises the covariance of
the projected factors with the responses,
%
\begin{equation}
\bm w_a = \arg\max_{\|\bm w\|=1}\;
\Cov\!\bigl(\bm X_s \bm w,\; \bm Y_s \bm c\bigr).
\end{equation}
%
Two diagnostics make PLS the standard tool of chemometrics and QbD. The first is
the \emph{Variable Importance in Projection}, which ranks how much each factor
contributes to the prediction,
%
\begin{equation}
\mathrm{VIP}_j = \sqrt{\,p \sum_{a}
  \frac{\mathrm{SS}_a}{\sum_b \mathrm{SS}_b}
  \left(\frac{w_{ja}}{\|\bm w_a\|}\right)^{\!2}},
\qquad
\mathrm{SS}_a = (\bm t_a\transp\bm t_a)(\bm q_a\transp\bm q_a),
\label{eq:vip}
\end{equation}
%
where $p$ is the number of factors; a factor with $\mathrm{VIP}_j>1$ is more
important than average. The second is the cross-validated coefficient of
prediction $Q^2$: the data are split into $k$ folds, each held out in turn and
predicted from a model fitted on the rest, and
%
\begin{equation}
Q^2 = 1 - \frac{\sum_{\text{folds}} \|\bm y_{\text{test}}
                 - \hat{\bm y}_{\text{test}}\|^2}
               {\sum \|\bm y - \bar{\bm y}\|^2}.
\end{equation}
%
Whereas the ordinary $R^2$ always rises as components are added, $Q^2$ measures
genuine \emph{predictive} power and falls when the model begins to fit noise,
giving an honest stopping rule for the number of latent variables.

\begin{prosbox}{PCA / PLS}
\begin{itemize}
  \item \textbf{Handle collinearity.} Both thrive where ordinary regression
  fails---many correlated factors, more variables than runs.
  \item \textbf{Insight, not just prediction.} Score and loading plots reveal
  structure; VIP ranks the drivers; $Q^2$ validates honestly.
  \item \textbf{Multi-response.} PLS models several CQAs jointly, exploiting their
  correlations.
\end{itemize}
\end{prosbox}

\begin{consbox}{PCA / PLS}
\begin{itemize}
  \item \textbf{Linear, latent.} Both assume the structure lives in linear
  combinations of the inputs; strong nonlinearity or interaction is captured only
  by adding engineered terms.
  \item \textbf{Scaling-dependent.} Results hinge on the (here automatic)
  pre-scaling; a poor choice distorts the components.
  \item \textbf{Indirect interpretation.} Latent variables are mathematical
  constructs, not physical knobs; the VIP and loadings need careful reading.
\end{itemize}
\end{consbox}

%==============================================================================
\section{Adaptive design: Bayesian optimisation}
\label{sec:bo}

Factorial and space-filling designs are committed in advance. When experiments
are expensive and the goal is to \emph{find the best} conditions rather than to
characterise the whole space, one can do far better by choosing each new
experiment in light of all previous ones. This is \emph{Bayesian optimisation}
(BO), and it is the package's most sophisticated design method
(\code{optimization/}).

\subsection{The Gaussian-process surrogate}

BO maintains a probabilistic \emph{surrogate} of the response: a Gaussian process
(GP). A GP places a distribution over functions such that any finite set of
points is jointly Gaussian, specified by a mean (taken zero after
standardisation) and a covariance \emph{kernel} $k(\bm x,\bm x')$. The package
uses \code{botorch}'s \code{SingleTaskGP} with the default Mat\'ern-$5/2$ kernel,
%
\begin{equation}
k(\bm x,\bm x') = \sigma_f^2
\left(1 + \sqrt5\,r + \tfrac{5}{3} r^2\right)
e^{-\sqrt5\, r},
\qquad
r = \sqrt{\textstyle\sum_d (x_d-x'_d)^2/\ell_d^2},
\end{equation}
%
smooth enough for a chromatographic response surface yet not as rigid as the
squared-exponential. Conditioned on the observed data
$(\bm X,\bm y)$, the posterior at a new point $\bm x_\star$ is Gaussian with the
closed-form mean and variance
%
\begin{align}
\mu(\bm x_\star) &= \bm k_\star\transp (\bm K + \sigma_n^2 \bm I)^{-1}\bm y, \\
s^2(\bm x_\star) &= k(\bm x_\star,\bm x_\star)
  - \bm k_\star\transp (\bm K + \sigma_n^2 \bm I)^{-1}\bm k_\star,
\end{align}
%
where $\bm K_{ij}=k(\bm x_i,\bm x_j)$ and $(\bm k_\star)_i = k(\bm x_i,\bm
x_\star)$. The crucial output is not just the prediction $\mu$ but the
\emph{uncertainty} $s$, which is large where data are sparse. The kernel
hyperparameters---the length scales $\ell_d$, signal variance $\sigma_f^2$, and
noise $\sigma_n^2$---are fitted by maximising the exact marginal log-likelihood
(\code{fit\_gpytorch\_mll}). The wrapper (\code{GPSurrogate}) normalises inputs to
the unit cube and standardises outputs, then un-transforms predictions, so the
caller works throughout in physical units.

\subsection{The acquisition function}

The surrogate is turned into a decision by an \emph{acquisition function} that
scores how worthwhile it is to sample a candidate point, balancing
\emph{exploitation} (sampling where the mean is high) against \emph{exploration}
(sampling where the uncertainty is high). The package uses \emph{Expected
Improvement} (EI): the expected amount by which a new observation would beat the
current best $f^\star$,
%
\begin{equation}
\mathrm{EI}(\bm x) = \E\!\bigl[\max(f(\bm x)-f^\star,\,0)\bigr]
= (\mu-f^\star)\,\Phi(z) + s\,\phi(z),
\qquad z = \frac{\mu-f^\star}{s},
\label{eq:ei}
\end{equation}
%
with $\Phi$ and $\phi$ the standard normal CDF and PDF. The first term rewards a
high mean, the second a high variance; their balance is automatic and
parameter-free. To avoid floating-point underflow when improvement is tiny, the
package optimises the numerically stable \code{LogExpectedImprovement} over the
domain (\code{optimize\_acqf}, multi-start L-BFGS-B), and evaluates the oracle at
the maximiser.

\subsection{The loop}

The complete procedure (\code{bayesian\_optimization}) is:
%
\begin{enumerate}
  \item \textbf{Seed} with $n_{\mathrm{initial}}$ Latin-hypercube points
  (Section~\ref{sec:lhs}) and evaluate the oracle at each.
  \item \textbf{Iterate}: fit the GP to all data; maximise EI to propose the next
  point; evaluate the oracle there; append and repeat.
\end{enumerate}
%
Constraints (for example, ``purity $\ge 0.95$'') are handled by a penalty: an
infeasible evaluation is recorded with its objective set to zero so it cannot
become the running best. The natural extension to true multi-objective
exploration of a yield/purity \emph{Pareto front} is to swap EI for the expected
hypervolume-improvement acquisition (qNEHVI); the surrogate machinery is
unchanged.

\subsection{Why it wins}

Figure~\ref{fig:bo} shows the payoff. Against a one-shot factorial spending its
runs on a fixed grid, the BO loop concentrates evaluations where the response is
promising and climbs to the optimum in far fewer experiments. The package
quantifies this with \emph{running-best} and \emph{regret} curves
(\code{compare\_to\_doe}), the regret being the gap between the running best and
the global best found by either method.

\begin{figure}[t]
\centering
\includegraphics[width=0.85\linewidth]{bo_convergence.png}
\caption[Bayesian optimisation versus a one-shot design]{Running-best objective
versus number of experiments for Bayesian optimisation (blue) and a one-shot
factorial design (red) on the same problem. After a short space-filling seed, BO
exploits its surrogate to reach the optimum with markedly fewer evaluations---the
advantage that matters most when each experiment is costly.}
\label{fig:bo}
\end{figure}

\begin{prosbox}{Bayesian optimisation}
\begin{itemize}
  \item \textbf{Sample-efficient.} Reaches good optima in a fraction of the
  experiments a fixed design needs---decisive when each run is expensive.
  \item \textbf{Principled exploration.} The GP's uncertainty drives a rational
  explore/exploit balance; no tuning of a trade-off parameter.
  \item \textbf{Built-in uncertainty.} The surrogate yields calibrated error
  bars, useful in their own right, and extends cleanly to constraints and
  multiple objectives.
\end{itemize}
\end{prosbox}

\begin{consbox}{Bayesian optimisation}
\begin{itemize}
  \item \textbf{Sequential.} Each step needs the previous result, so it cannot be
  fully parallelised without batch acquisitions, and it ill suits campaigns where
  experiments must be committed in blocks.
  \item \textbf{Surrogate assumptions.} Performance rests on the GP being a fair
  model; deceptive, discontinuous, or very high-dimensional landscapes degrade
  it, and GP fitting scales cubically in the number of points.
  \item \textbf{Optimum, not map.} It finds the best point, not an interpretable
  effects table or a fully characterised design space---the deliverable a
  regulatory filing expects.
  \item \textbf{Less transparent.} The decision logic is harder to explain to
  reviewers than a factorial ANOVA.
\end{itemize}
\end{consbox}

%==============================================================================
\section{Choosing among the methods}
\label{sec:design-compare}

These families are complements, not rivals, and a mature characterisation
campaign uses them in sequence (Table~\ref{tab:compare}). One \emph{screens} many
candidate factors with a (fractional) factorial to find the vital few; one
\emph{explores} the reduced space with an LHS and \emph{understands} it with
PCA/PLS; one \emph{optimises} within it with BO; and one \emph{characterises} the
final design space with a response-surface design---and, where the deliverable is
the probability of meeting specification, a logistic GLM
(Section~\ref{sec:glm})---for the filing. The same
mechanistic oracle of Chapter~\ref{ch:models} drives them all, which is precisely
the point of a virtual laboratory: it lets one rehearse this entire strategy,
and quantify its uncertainty, before committing a single millilitre of resin.

\begin{table}[h]
\centering
\renewcommand{\arraystretch}{1.35}
\caption{The experimental-design methods at a glance.}
\label{tab:compare}
\begin{tabular}{@{}p{0.20\linewidth} p{0.26\linewidth} p{0.22\linewidth}
                  p{0.22\linewidth}@{}}
\toprule
\textbf{Method} & \textbf{Primary purpose} & \textbf{Adaptive?} &
\textbf{Best when} \\
\midrule
Full factorial / RSM & Estimate effects, interactions, design space &
No (one-shot) & Few factors; need interpretable, defensible effects \\
GLM (logistic / Poisson) & Probability of meeting spec; counts &
Analysis step & Pass/fail or count responses; probabilistic design space \\
Latin hypercube & Space-filling exploration; surrogate training data &
No (one-shot) & Building a global model; moderate-to-many factors \\
PCA / PLS & Understand structure; rank factors; relate $X\!\to\!Y$ &
Analysis step & Many correlated factors and responses \\
Bayesian optimisation & Find the optimum in fewest runs &
Yes (sequential) & Experiments are expensive; goal is an optimum \\
\bottomrule
\end{tabular}
\end{table}

\begin{keybox}{The throughline}
A faithful model (Chapter~\ref{ch:models}) and an efficient way to query it
(Chapter~\ref{ch:design}) are two halves of one instrument. The model supplies
\emph{physics}---the right shape and sensitivities of the response; the designs
supply \emph{strategy}---which questions to ask and how to learn from the
answers. Together they turn process development from a sequence of expensive
guesses into a directed, quantifiable search.
\end{keybox}
